%        File: Bball.tex
%     Created: Tue May 29 03:00 PM 2012 C
% Last Change: Tue May 29 03:00 PM 2012 C
%
\documentclass[12pt]{amsart}
\usepackage{latexsym,amssymb}
\usepackage{graphics}
\usepackage{url}
\usepackage{harvard}
\usepackage{eulervm}
\usepackage{epsfig,subfigure}
\usepackage{color}

\setlength{\textheight}{205mm}
\setlength{\oddsidemargin}{10mm}
\setlength{\evensidemargin}{10mm}
\setlength{\topmargin}{0mm}
\setlength{\textwidth}{6in}
\def\RR{\mathbb{R}}
\def\EB{\mathbb{E}}
\def\VB{\mathbb{V}}
\def\OO{\mathcal O}
\def\DD{\mathcal D}
\def\BB{\mathcal B}
\def\PP{\mathbb{P}}
\def\LL{\mathcal L}
\def\II{\mathcal I}
\def\EE{\mathcal E}
\def\FF{\mathcal F}
\def\NN{\mathcal N}
\def\XX{\mathcal X}
\def\SS{\mathcal S}
\def\GG{\mathcal G}
\def\CC{\mathcal C}
\def\AA{\mathcal A}
\def\TT{\mathcal T}
\def\HH{\mathcal H}
\def\JJ{\mathcal J}
\def\KK{\mathcal K}
\def\PP{\mathbb{P}}
\def\LL{\mathcal L}
\def\II{\mathcal I}

\def\Polya{P\'olya }
\def\Renyi{R\'enyi }
\def\diag{\mbox{diag}}
\def\TV{\bigvee_\Omega }
\newtheorem{example}{Example}
\newtheorem{remark}{Remark}
\renewcommand{\abstractname}{Abstract}
\def\argmax{\text{argmax}}
\newtheorem{lemma}{Lemma}
\newtheorem{theorem}{Theorem}
\newtheorem{corollary}{Corollary}
\newtheorem{proposition}{Proposition}
\renewenvironment{proof}{\noindent {\bf Proof.\ }}{\hfill{\rule{2mm}{2mm}}}
\renewenvironment{remark}{\noindent {\bf Remark.\ }} {\hfill{ \rule{2mm}{2mm}}}
\renewenvironment{example}{\noindent {\bf Example.\O}}


\begin{document}
\bibliographystyle{econometrica}
\title{Unobserved Heterogeneity in Longitudinal Data: \\
An Empirical Bayes Perspective}

\author{Jiaying Gu}
\author{Roger Koenker}
\thanks{Preliminary Version:  \today .    This research was partially supported by NSF grant
SES-11-53548.  All computational results of this paper can be reproduced with
the R packages REBayes, \citeasnoun{REBayes},  and Rmosek, \citeasnoun{Rmosek}, 
and code available from the authors on request.
The Rmosek package provides a convenient R interface to the Mosek optimization
language of \citeasnoun{mosek}.  We would like to express our appreciation to
participants in the CeMMaP-UCL, Harvard-MIT, Georgetown, Rutgers, Cornell and 
University of Maryland econometrics
workshops for valuable comments while retaining full responsibility for all errors of
omission and commission.}

\begin{abstract}
	Empirical Bayes methods for Gaussian and binomial compound decision problems
	involving longitudinal data are considered.  A new convex optimization formulation of 
	the nonparametric (Kiefer-Wolfowitz) maximum likelihood estimator for mixture models is used
	to construct nonparametric Bayes rules for compound decisions.  The methods
	are illustrated with some simulation examples as well as an application to
	predicting baseball batting averages.  Comparisons with nonparametric Bayesian
	methods based on Dirichlet process priors are also provided.
\end{abstract}
\maketitle
\pagestyle{myheadings}
\markboth{\sc Empirical Bayesball}{\sc Gu and Koenker}
%\begin{quote}
%	{\footnotesize{\it
%	It is a truth universally acknowledged, that a statistician in
%	possession of a good shrinkage estimator must be desirous of 
%	predicting baseball batting averages. 
%	
%	\hspace{50mm} Jane (Yogi) Austen}}
%	%It is a truth universally acknowledged, that a single man in 
%	%possession of a good fortune, must be in want of a wife. (Jane Austen)
%\end{quote}
\section{Introduction}
Unobserved heterogeneity has become a pervasive concern throughout applied econometrics.  
%No where is this more true than in survival analysis where 
%unobserved heterogeneity is often referred to as frailty.  
Longitudinal data presents special opportunities and challenges for models of 
unobserved heterogeneity; in virtually all econometric applications involving
panel data there will be some form of latent, i.e. unobserved, individual specific
effects.  Classical methods adopt either a differencing strategy designed to
purge these effects, or some form of shrinkage method to mitigate their undesirable
``incidental parameter'' effect.  In this paper we will describe some new 
nonparametric empirical Bayes methods for estimation and
prediction in panel data models with unobserved heterogeneity, and compare
and contrast them with some existing nonparametric Bayesian proposals.

As stressed in recent work of Efron (2010, 2011), empirical Bayes methods pioneered
\nocite{Efron.10}
\nocite{efron.11}
by Robbins (1951, 1956) provide a statistical framework for many contemporary
``big data''  applications.  Although they predate the development of hierarchical
Bayes methods exemplified in the work of \citeasnoun{LS} and \citeasnoun{LC},
they share many common features.  The transition from parametric to nonparametric
empirical Bayes methods brings exciting new opportunities that greatly expand the
flexibility of existing approaches to panel data modeling and its treatment of
unobserved heterogeneity.

We will begin with a brief review of some developments in empirical Bayes methods
beginning with \citeasnoun{Robbins.51}, touching on the connections to Stein rule 
methods and finally evolving into modern nonparametric mixture variants.  In
Section 3 we extend the predominant Gaussian location mixture framework to accommodate
nonparametric location {\it and scale} mixtures in the classical Gaussian panel
data setting.  Section 4 presents some simulation evidence to illustrate the
performance of the new methods, and Section 5 describes an extended application
that illustrates both estimation and prediction aspects of the new methods including,
notably, the introduction of covariate effects via profile likelihood methods.

In sharp contrast to the classical Gaussian hierarchical Bayes framework for panel
data, or its frequentist analogues, the nonparametric mixture formulation of our
proposed methods offers a much more flexible view of unobserved heterogeneity
while preserving most of the virtues of the likelihood formalism. 

\section{Empirical Bayes:  A Brief Review}

Given a simple parametric model, there is a natural temptation
to complicate it by admitting that those immutable natural constants that
constitute the model's original parameters might instead be random.  One of the
earliest examples of this type is the classical Gaussian random effects,
compound decision problem introduced by \citeasnoun{Robbins.51}.  We observe
independent $Y_1, \cdots , Y_n$ each Gaussian with common variance, $\sigma^2$
but individual specific means, $Y_i \sim \NN (\mu_i , \sigma^2)$.  Our 
objective is to estimate all the $\mu_i$'s subject to squared error loss,
\[
\LL_2 (\hat \mu , \mu ) = \| \hat \mu - \mu \|_2^2 = 
\sum_{i=1}^n (\hat \mu_i - \mu_i)^2.
\]
The naive (unbiased) solution would simply set $\hat \mu_i = Y_i$,
but the usual presumption in such circumstances would be that the observations
have some common genesis, and consequently we may be able to ``borrow
strength''  from the full sample to improve upon these myopic predictions
based on a single observation.

Suppose we believed that the $\mu_i$ were drawn
iid-ly from the distribution, $F$, so the $Y_i$'s would have convolution
density $g(y) = \int \phi ((y - \mu)/\sigma)/\sigma dF(\mu)$:  What would
the Reverend Bayes advise?  Elementary exponential family theory yields
the following proposition.  
%Appendix A contains a concise proof. 

\begin{proposition}
	For $Y_i \sim \NN(\mu_i , \sigma^2)$ and $\{ \mu_i \}$ iid $F$,
	the Bayes rule under $\LL_2$ loss is:
	\begin{equation}
	\delta(y) = y + \sigma^2 g^\prime (y) / g(y)
	\label{eq.Tweedie}
	\end{equation}
	and $\delta (y)$ is non-decreasing in $y$.
\end{proposition}
\citeasnoun{efron.11} refers to  \eqref{eq.Tweedie} as Tweedie's formula
citing Robbins's (1956) \nocite{Robbins.56} attribution of it to M.C.K.
Tweedie.  Of course one may well ask:  Where did this $F$ come from?  And
this question leads us inevitably toward estimation of the density, $g$, and 
hence to the empirical Bayes paradigm.  When $F$ comes from a finite 
dimensional parametric family there are several prominent special cases.

\subsection{Some Parametric Examples}
\begin{enumerate}
	\item Suppose $\sigma^2 = 1$ and we believed that the $\mu_i$'s were iid
$\NN (0 , \sigma_0^2)$, so the $Y_i$'s are iid $\NN (0 , 1 + \sigma_0^2)$, 
the Bayes rule would be,
\[
\delta (y) = \left ( 1 - \frac{1}{1 + \sigma_0^2} \right ) y.
\]
Thus, we shrink our naive estimates all toward zero.  When $\sigma_0^2$
is unknown, $S = \sum Y_i^2 \sim (1 + \sigma_0^2) \chi_n^2$, and
recalling that an inverse $\chi_n^2$ random variable has expectation,
$(n-2)^{-1}$, we obtain the Stein rule in its simplest form:
\[
\hat \delta (y) = \left ( 1 - \frac{n-2}{S} \right ) y.
\]
\item When, slightly more generally, $\mu_i \sim \NN (\mu_0 , \sigma_0^2)$
we shrink instead  toward the prior mean,
\[
\delta (y) = \mu_0 + \left ( 1 - \frac{1}{1 + \sigma_0^2} \right ) (y - \mu_0),
\]
and estimating the prior parameters costs us one degree of freedom, so
we obtain the celebrated James-Stein estimator,
\[
\hat \delta (y) = \bar Y_n + \left ( 1 - \frac{n-3}{S} \right ) 
(y - \bar Y_n),
\]
for $\bar Y_n = n^{-1} \sum Y_i$ and $S = \sum (Y_i - \bar Y_n)^2.$

\item If each observation has its own {\it known} variance:
$Y_i \sim \NN(\mu_i, \sigma_i^2)$ and $\mu_i \sim \NN(\mu_0 , \sigma_0^2)$,
as might be plausible in the case that each $Y_i$ is from a different 
measuring device each with known precision, or as in the ubiquitous baseball 
batting average examples, as in \citeasnoun{brown.08} and \citeasnoun{JZ.10},
in which binomial variances depend upon a known number of ``at bats'' in 
the initial period.  In such cases we have the Bayes rule,
\[
\delta (y_i) 
= \mu_0 + \left( 1 - \frac{\sigma_i^2}{\sigma_0^2 + \sigma_i^2} \right ) (y_i - \mu_0 )
\]
\item Further generalizing, we may wish to replace $\mu_0$ by a function of 
observable covariates, say $z_i^\top \beta_0$. Then, as in
\citeasnoun{JZ.10}, we obtain a positive-part James-Stein estimator,
\[
\hat \delta (y_i) = \left( 1 - 
\frac{p-2}{\sum (z_i^\top \hat \beta/\sigma_i)^2} \right )_+ 
z_i^\top \hat \beta + 
\left (1 - \frac{n - p-2}{\sum (y_i - z_i^\top \hat \beta)^2/\sigma_i^2} \right )_+ 
(y_i - z_i^\top \hat \beta)
\]
where $p$ denotes the dimension of $\beta$ and $(u)_+ = u I(u > 0)$.  

\item Another important class of examples arises from the assumption  of sparsity,
	that is, an assertion that most of the $\mu_i$ are probably zero.  \citeasnoun{JS}
	consider a model in which,
	\[
	dF(\mu) = (1-w) \delta_0 (\mu) + w \varphi_\nu (\mu )
	\]
	where with probability $(1-w)$, $\mu = 0$, while with probability $w$
	it is drawn from a density, $\varphi$, with scale parameter, $\nu$.
	They compare performance of several hard and soft threshholding rules
	in addition to empirical Bayes procedures that estimate the parameters
	$w$ and $\nu$.  
\end{enumerate}
The simulation designs of \citeasnoun{JS} have served as a benchmark
for several more recent studies of empirical Bayes methods including
\citeasnoun{BG}, \citeasnoun{JZ}, and  \citeasnoun{KM},  all of which
explore non-parametric estimation of the Gaussian mixture model.

\subsection{Non-parametric Estimation of the Gaussian Mixture Model}
When we lack confidence in a particular parametric specification of
the mixing distribution, $F$, we are faced with a serious quandary.
It is apparent that we need a non-parametric estimate of $F$, and in
our Gaussian location mixture setting this is tantamount to solving
a deconvolution problem.  Deconvolution is notoriously difficult as demonstrated
by \citeasnoun{carrollhall.88} and \citeasnoun{fan.91}, but before we
despair a second look at Tweedie's formula \eqref{eq.Tweedie} reveals
that we don't really {\it need} an estimate of $F$, we need only estimate
the mixture density $g$, a task that can be accomplished at standard univariate
non-parametric convergence rates for smooth densities.  Smoothness is ensured
by the Gaussian convolution whatever $F$ might be.

Kernel density estimation as proposed by \citeasnoun{BG} seems to be the natural
approach, but in addition to the familiar, but still unsettling, requirement 
of selecting a bandwidth, the kernel estimator has the drawback that 
it does not enforce the monotonicity of the Bayes rule.  The
latter failing can be addressed by a further monotonization step, or by a
penalization approach as suggested in \citeasnoun{KM}.  However, a more direct
approach is possible via the Kiefer-Wolfowitz non-parametric maximum likelihood
estimator.  This approach was first proposed by \citeasnoun{JZ} for the
Gaussian compound decision problem, employing the EM algorithm as a computational
strategy. 

The EM algorithm has been the standard approach to computing the 
Kiefer-Wolfowitz non-parametric MLE since it was originally suggested by 
\citeasnoun{Laird}; \citeasnoun{HeckmanSinger}
constitutes an influential econometric application.  However, EM has notoriously slow
convergence in these applications, and this fact introduces what is, in effect, a
new smoothing parameter into the computational strategy controlled by the stopping
criterion of the algorithm.  \citeasnoun{KM}  have recently proposed an alternative
computational method for the Kiefer-Wolfowitz MLE that circumvents these problems.  For
a broad class of mixture problems, the Kiefer-Wolfowitz estimator can be formulated 
as a convex optimization problem and solved efficiently by modern interior point methods.  

In the next section we will describe how these methods can be extended 
to longitudinal data, first for location and scale mixtures
separately, then for location-scale mixtures and finally for location
scale mixtures with covariate effects.  In contrast to compound 
decision problems with cross sectional data, richer longitudinal data 
offers new opportunities permitting more complex structure of  unobserved 
heterogeneity. 

\section{Estimating Gaussian Mixtures with Longitudinal Data}

Extending the Gaussian compound decision problem with one location parameter per
observation to unbalanced longitudinal observations in which we have $m_i$ 
observations on each individual is straightforward.  Assuming for convenience
that we have unit variance for the noise so $u_{it} \sim \NN(0,1)$, we have,
\[
y_{it} = \mu_i + u_{it}, \quad t = 1, \cdots , m_i, \quad i = 1, \cdots , n,
\]
but sufficiency reduces the problem to our third example with 
$\hat \mu_i = m_i^{-1} \sum_{i=1}^{m_i} y_{it} \sim \NN (\mu_i , m_i^{-1} )$.
When the $\mu_i$'s are iid from $F$, we can write the log likelihood of the
observed $y_{it}$'s as,
\[
\ell (F | y ) = \sum_{i=1}^n \log ( \sqrt{m_i} \int 
\phi( \sqrt{m_i} (\hat \mu_i - \mu)) dF(\mu)
\]
Optimizing over the infinite dimensional $F$ necessitates some form of discrete
approximation.  As in earlier EM implementations, such as that of \citeasnoun{JZ},
we take $F$ to have a piecewise constant (Lebesgue) density on a relatively fine
grid containing the empirical support of the observed $\hat \mu_i$'s.  With
a few hundred grid intervals we can obtain a quite accurate estimate.  Further
refinement is always possible as discussed in \citeasnoun{KM}, but already with
a uniform grid with 300 points we have very precise positioning of the mass points
of the mixing distribution, more precise than the statistical accuracy of the mass
locations would justify.  Letting $f_j :
j = 1, \cdots , p$ denote the function values of $dF$ on this grid, we can express
the constrained maximum likelihood problem as,
\begin{equation}
	\max_f \{ \sum_{i=1}^n \log (g_i) \; | \; g = A f , \; 
	\sum_{j=1}^p f_j \Delta_j = 1, \; f \geq 0 \},
	\label{eq.PL}
\end{equation}
where $A = (A_{ij} = ( \sqrt{m_i} \int \phi( \sqrt{m_i} (\hat \mu_i - \mu)))$ and
$\Delta_j$ is the $j$th grid spacing.  As posed, the problem is evidently
convex, and therefore has a unique solution.  It is well-known, going back to
\citeasnoun{KW} and \citeasnoun{Laird}, that variational solutions to the
original problem are discrete with fewer than $n$ atoms.  It is somewhat
difficult to appreciate this result by viewing EM solutions, since the number
of EM iterations required to obtain an accurate solution would test the patience
of the most diligent researchers.  But interior point methods make this discreteness
easily apparent.  Since the number of non-negligible $\hat f_i > 0$ obtained is
typically much smaller than $n$, often only a handful of points, even in 
large samples, this also guides our judgement regarding the adequacy of the original grid.
Again, larger $n$ might justify a refinement of the grid at very modest increase in
computational effort.

The dual formulation of primal problem \eqref{eq.PL} has proven to be somewhat more 
efficient from a computational viewpoint.  The dual can be expressed as
\begin{equation}
	\max_\nu \{ \sum_{i=1}^n \log (\nu_i) |  A^\top \nu \leq n 1_n, \nu \geq 0 \},
	\label{eq.DL}
\end{equation}
see \citeasnoun{KM} for further details.




\subsection{Estimating Gaussian Scale Mixtures}
Gaussian scale mixtures can be estimated in much the same way that we have
described for location mixtures.  Suppose we now observe,
\[
y_{it} = \sigma_i u_{it}, \quad t = 1, \cdots , m_i, \quad i= 1, \cdots , n
\]
with $u_{it} \sim \NN (0,1)$.  Sufficiency again reduces the sample to
$n$ observations on $s_i = m_i^{-1} \sum_{t=1}^{m_i} y_{it}^2$, and thus
$2 r_i s_i/\sigma_i^2$ with $r_i = m_i / 2$ has the gamma distribution with
shape parameter, $r_i$ and scale parameter $\sigma_i^2 / r_i$, i.e.
\[
\gamma (s_i | r_i,  \sigma_i^2 ) = \frac{1}{\Gamma (r_i) (\sigma_i^2/r_i)^{r_i}}
s_i^{r_i-1} \exp \{ - s_i r_i /\sigma_i^2 \},
\]
and the marginal density of $s_i$ when the $\sigma_i^2$ are iid from $F$ is
\[
g(s_i ) = \int \gamma (s_i | r_i, \sigma^2 ) dF(\sigma^2).
\]
To estimate $F$ we can proceed exactly as before except that now the matrix
$A$ has typical element $\gamma(s_i | \sigma_j^2)$ for $\sigma_j^2$ on
a fine grid covering the support of the sample $s_i$'s. 
\subsection{Estimating Gaussian Location-Scale Mixtures}
When both location and scale are heterogeneous we must combine the strategies
already described.  We should emphasize here that the scope for modeling 
heterogeneity for the scale parameter would not be possible with cross 
sectional data since individuals are then  only measured once. The model is now,
\[
y_{it} = \mu_i  + \sigma_i u_{it}, \quad t = 1, \cdots , m_i, 
\quad i= 1, \cdots , n
\]
with $u_{it} \sim \NN (0,1)$.  We will provisionally assume that $\mu_i \sim F_\mu$ and
$\sigma_i^2 \sim F_{\sigma^2}$ are independent.  Again, we have sufficient
statistics:
\[
\hat \mu_i | (\mu_i , \sigma_i^2) \sim \NN (\mu_i, \sigma_i^2/m_i )
\]
and
\[
s_i | (r_i , \sigma_i^2 ) \sim \gamma ( s_i | r_i , \sigma_i^2/r_i ),
\]
where $r_i = ( m_i - 1) /2$, and the log likelihood becomes,
\[
\ell (F_\mu , F_{\sigma^2} | y) = \sum_{i=1}^n \log \int \int
\gamma(s_i | r_i , \sigma^2 /r_i ) \sqrt{m_i} 
\phi (\sqrt{m_i} (\hat \mu_i - \mu_i)/\sigma)/\sigma 
dF_\mu (\mu) dF_{\sigma^2} (\sigma^2)
\]
Since the scale component of the log likelihood is additively separable from the location
component, we can solve for $\hat F_{\sigma^2}$ in a preliminary step, as
in the previous subsection, and then solve for the $\hat F_\mu$ distribution.
In fact, under the independent prior assumption, we can  re-express the Gaussian
component of the likelihood as Student-$t$ and thereby eliminate the dependence on
$\sigma$ in the Kiefer-Wolfowitz problem for estimating $F_\mu$.   This is highly
convenient for estimation purposes, however it should be stressed that prediction
restores the interdependence on both $F_\mu$ and $F_\sigma$ as we discuss in more
detail below.

When the independence assumption is implausible, and this may be typical of econometric
applications where there is some portfolio aspect of the problem that makes $\mu_i$'s and
$\sigma_i$'s positively correlated, we can construct two dimensional grids.  This makes
the constraint matrix, $A$, somewhat unwieldy, but raises no new issues in principle.
After a brief discussion of testing for parameter heterogeneity, we will compare and
contrast the foregoing methods in simulations and an application to predicting 
baseball batting averages.



A natural question at this point might be:  How do we know whether we {\it need}
to bother with all this?  Can we make some preliminary test for parameter
heterogeneity?  There is an extensive literature on this topic:
\citeasnoun{Chesher.84} and \citeasnoun{Cox.83} constitute notable contributions  
drawing connections to the \citeasnoun{White82} information matrix test.
Many of the proposals that have been made can be formulated
as $C(\alpha)$ tests as in \citeasnoun{Neyman.59} and  
\citeasnoun{NeymanScott.66}.  The $C(\alpha)$ formulation can be viewed
as an extended form of the Rao score test that among other things can
accommodate general forms of nuisance parameter estimation. 
\citeasnoun{Gu13} provides a detailed theoretical treatment of
$C(\alpha)$ tests for parameter heterogeneity, and  \citeasnoun{GKV} compares 
their performance to that of likelihood ratio tests.
In Appendix B we provide a detailed discussion of the $C(\alpha)$ test for homogeneity
for variances in the longitudinal Gamma model that is employed in the application section.


\section{Empirical Bayes Prediction:  Some Simulation Evidence}
To develop some intuition about empirical Bayes methods we will consider some 
simple illustrative examples in this section before turning to our main
empirical application.

\subsection{Gaussian Location Mixtures}
Suppose that we have a random sample from the model: $y_i = \mu_i + u_i$ with
iid $u_i \sim \NN (0,1)$, and iid $\mu_i \sim \frac{2}{3} \delta_{-h} + 
\frac{1}{3} \delta_{2h}$ as in \citeasnoun{Chen.95}.  Here, $\delta_a$ denotes
the distribution with point mass one at the point $a$.  If we were successful
in estimating the $\mu_i$ distribution, we would expect that Tweedie's formula
\eqref{eq.Tweedie} should deliver predictions that correctly shrink the
original observations toward their respective $\mu_i$'s.  In this sense the
prediction problem becomes one closely related to classification.  In Figure \ref{fig.ex1}
we illustrate this situation with $n = 400$ and $h = 0.5$, which represents a 
fairly challenging problem.  In the left panel we have the estimated mixing
distribution in red, with the target distribution represented in blue.
The larger of the two actual mass points at $x = -0.5$ is quite accurately
estimated, however the smaller mass point at $x = 1$ is split into two
pieces by the Kiefer-Wolfowitz estimate.  The true mixture distribution
in blue in the middle panel appears to be reasonably accurately estimated
by the red curve.  The corresponding Bayes rules in the right panel show
that the empirical Bayes rule (in red) shrinks a little too aggressively
in the left tail, and not quite agressively enough in the right tail, compared
to the omniscient  Bayes rule in blue.  But we should hasten to add that it
represents an enormous improvement over the unbiased (naive) decision rule,
$\hat \mu_i = y_i$, depicted in grey.


\begin{figure}[ht]
  \begin{center}
          \resizebox{ \textwidth}{!}{{\includegraphics{Tweedie/ex1.pdf}}}
  \end{center}
  \caption{{Empirical Bayes estimation for the Chen (1995) example:  A sample
  of 400 observations from the model with $y_i = \mu_i + u_i$ with
  iid $u_i \sim \NN (0,1)$, and iid $\mu_i \sim \frac{2}{3} \delta_{-h} + 
  \frac{1}{3} \delta_{2h}$, is illustrated by the histogram in the middle
  panel and the ``rug plots'' in the adjacent panels.  The Kiefer-Wolfowitz
  estimate of the mixing distribution is illustrated in the left panel in
  red, with the target distribution in blue.  The corresponding estimate of
  the mixture density and Bayes rule appear in the other panels contrasted 
  to their blue target functions.  The unbiased, naive decision rule is
  depicted in the right panel in grey.}}
  \label{fig.ex1}
\end{figure}

Replacing the two mass point distribution by $\mu_i \sim U[-h, 2h]$
yields an even more challenging problem.  The Kiefer-Wolfowitz  estimator
tries valiantly to mimic the uniform mixture by a discrete mixture as
illustrated in Figure \ref{fig.ex2}.  The two point mixing distribution 
appearing in the left panel does not seem to be a very satisfactory 
surrogate for the uniform, but as can be seen in the middle panel, it
does a remarkably good job of imitating the correct mixture density.
The Bayes rule comparison in the right panel again illustrates that
the shrinkage in the tails is not ideal, but much preferable to the
naive, unbiased rule.


\begin{figure}[ht]
  \begin{center}
          \resizebox{ \textwidth}{!}{{\includegraphics{Tweedie/ex2.pdf}}}
  \end{center}
  \caption{{Empirical Bayes estimation for the Chen (1995) example:  A sample
  of 400 observations from the model with $y_i = \mu_i + u_i$ with
  iid $u_i \sim \NN (0,1)$, iid $\mu_i \sim U[-h, 2h]$, and $h = 0.5$,
  is illustrated by the histogram in the middle
  panel and the ``rug plots'' in the adjacent panels.  The Kiefer-Wolfowitz
  estimate of the mixing distribution is illustrated in the left panel in
  red, with the target distribution in blue.  The corresponding estimate of
  the mixture density and Bayes rule appear in the other panels contrasted 
  to their blue target functions.  The unbiased, naive decision rule is
  depicted in the right panel in grey.}}
  \label{fig.ex2}
\end{figure}

\subsection{Gamma Scale Mixtures}

To explore the performance of empirical Bayes methods for gamma scale mixtures
we illustrate a couple of similar cases to those appearing in the previous
subsection.  We first consider the longitudinal model, 
\[
y_{it} = \mu_i + \sigma_i u_{it}, \quad t = 1, \cdots, m, \quad i=1,\cdots , n,
\]
with iid $u_{it} \sim \NN (0,1)$, and
$\sigma_i^2 \sim F$.
We take
$m = 11$ and $n = 400$.  We will provisionally ignore the heterogeneity
in the $\mu_i$, or to be more explicit, adopt the naive practice of estimating
them by $\bar y_i$.  Denoting the individual specific variance estimates by
$x_i =  (m-1)^{-1} \sum (y_{it} - \bar y_i)^2$, and their population
counterparts by $\theta_i = \sigma_i^2$, the $\{x_i\}$ are then distributed
as Gamma with shape parameter, $r = (m-1)/2$, scale parameter $\theta_i /r$,
and density,
\[
\gamma( x_i | \theta_i ) = \frac{1}{\Gamma(r) (\theta_i /r)^r} x_i^{r-1} 
\exp( - x_i r / \theta_i ).
\]
Thus, the marginal density of the sample variances is,
\[
g(x) = \int \gamma (x | \theta ) dF(\theta).
\]
The Bayes rule under squared error loss for $\theta_i$ given $x_i$,
originally derived by \citeasnoun{Robbins.82}, is given in the following 
proposition.  Again, it should be stressed that the Bayes rule depends only 
on the mixture density, $g$, and not directly on the mixing distribution, $F$.
Of course, indirectly the Bayes rule {\it does} depend on $F$ and in particular
the flat portions of the Bayes rule figures representing the points of attraction
of Bayes shrinkage are essentially determined by the location of the estimated mass
points of $F$.  This is particularly crucial in the tails where even small mass points
of $\hat F$ can exert a large influence on the shrinkage.  Whether this sensitivity
can be lessened by replacing the Gaussian mixture assumption by something heavier
tailed constitutes an intriguing question for future research.  The next proposition
describes the Bayes rule for Gamma mixtures.
\begin{proposition}
	For $X_i \sim \Gamma(r, \theta_i/r)$ and $\{ \theta_i \}$ iid $F$,
	the Bayes rule under $\LL_2$ loss is:
	\begin{equation}
		\delta(x) = r x^{r-1} \int_x^\infty y^{1-r} g(y) dy /g(x) 
	\label{eq.Tweediev}
	\end{equation}
	and $\delta (x)$ is non-decreasing in $x$.
\end{proposition}

\begin{figure}[ht]
  \begin{center}
          \resizebox{ \textwidth}{!}{{\includegraphics{Tweedie/exv1.pdf}}}
  \end{center}
  \caption{{Empirical Bayes estimation for Gamma mixture example:  A sample
  of $n =400$ and $m = 11$ observations from the model 
  $y_{it} = \mu_i + \sigma_i u_{it}$ 
  with iid $u_{it} \sim \NN (0,1)$, iid $\sigma_i^2 \sim
  \frac{2}{3} \delta_{1.5} + \frac{1}{3} \delta_3$.  
  is illustrated by the histogram in the middle
  panel and the ``rug plots'' in the adjacent panels.  The Kiefer-Wolfowitz
  estimate of the mixing distribution is illustrated in the left panel in
  red, with the target distribution in blue.  The corresponding estimate of
  the mixture density and Bayes rule appear in the other panels contrasted 
  to their blue target functions.  The unbiased, naive decision rule is
  depicted in the right panel in grey.  The brown line represents the 
  linearized empirical Bayes rule proposed in Robbins (1982)}}.
  \label{fig.exv1}
\end{figure}

In Figure \ref{fig.exv1} we illustrate a
typical outcome in a format like that of the previous  figures.  
In this example we take $F$ to be the two point distribution:
$\frac{2}{3} \delta_{1.5} + \frac{1}{3} \delta_3$.  
The two point mixing distribution is quite well estimated by the Kiefer-Wolfowitz
procedure, and the mixture density appears to be quite accurate as well.
The empirical Bayes rule slightly over estimates the variances in the upper
tail since it slightly overestimates the location of the upper mass point.
But as for the previous examples, there is an enormous improvement over the
naive decision rule represented by the grey line.  The brown line represents
the linearized empirical Bayes rule proposed in \citeasnoun{Robbins.82}.

\subsection{Gaussian Location Scale Mixtures}

We now would like to consider joint estimation of location and scale mixtures
in the context of our longitudinal model.  We will maintain the assumption that
$\mu_i$'s and $\sigma_i^2$'s are drawn independently, so we only have to estimate
two univariate mixing densities rather than a general bivariate density.
We illustrate the procedure with an example that combines a three point distribution
for $\mu$ and a three point distribution for $\sigma^2$.  Maximizing the likelihood
appear in Section 2.3, we obtain the estimates appearing in the first two panels
of Figure \ref{fig1.exv3} for the location and scale parameters respectively.  As
in the previous figures, estimates appear in red, and the true mixing distribution
is represented by the blue lines.  Focusing initially on the location parameter,
the third panel of the figure depicts the histogram of the observed $\hat \mu_i =
\bar y_i$ with the estimated marginal density superimposed in red, and the true
marginal density superimposed in blue.  Finally, in the
last panel of the figure we illustrate the empirical and idealized Bayes rule
for the $\mu_i$'s.  This version of the Bayes rule presumes that prediction is
based on knowledge of a location estimate, but nothing about the scale
parameter beyond the distribution represented by the estimated mixing distribution.
Even though the mixing distribution of the location parameter has a few extraneous
mass points, the Bayes rule is remarkably accurate.  

\begin{figure}[ht]
  \begin{center}
          \resizebox{ \textwidth}{!}{{\includegraphics{Tweedie/fig1exv3.pdf}}}
  \end{center}
  \caption{{Empirical Bayes estimation for Gaussian location-scale mixture:  
  A sample of $n = 800$ and $m = 11$ observations from the model 
  $y_{it} = \mu_i + \sigma_i u_{it}$ 
  with iid $u_{it} \sim \NN (0,1)$, 
  iid $\mu_i \sim \frac{1}{3} \delta_{-0.5} + \frac{1}{3} \delta_1+ \frac{1}{3} \delta_3$, and  
  iid $\sigma_i^2 \sim \frac{1}{3} \delta_{0.5} + \frac{1}{3} \delta_2 + \frac{1}{3} \delta_4$,  
  is illustrated by the histogram in the middle panel The Kiefer-Wolfowitz
  estimate of the mixing distributions is illustrated in the two left panels in
  red, with the target distribution in blue.  The corresponding estimate of
  the mixture density and the Bayes rule appear in the other panels contrasted 
  to their blue target functions.}} 
  \label{fig1.exv3}
\end{figure}

How much can be gained by using an individual specific estimate of variance?  To explore
this question,  we generate data from a simple Gaussian mixture model that
combines the prior three point mean and three point variance mixture:
$y_{it} = \mu_i + \sigma_i u_{it}$ with iid $u_{it} \sim \NN (0,1)$, 
iid $\mu_i \sim \frac{1}{3} \delta_{-0.5} + \frac{1}{3} \delta_1+ \frac{1}{3} \delta_3$, and  
iid $\sigma_i^2 \sim \frac{1}{3} \delta_{0.5} + \frac{1}{3} \delta_2 + \frac{1}{3} \delta_4$,  
A typical realization of this model is illustrated in Figure \ref{fig1.exv3}
where we see that both the mean and variance mixture are quite well estimated.
The marginal density of the $\bar y$'s is shown in the third panel, first as
simply a histogram of the observed values, then (in red) as an estimate based on the
Kiefer-Wolfowitz estimate of the mixing densities, and finally (in blue) as the
true population mixture.  In the far right panel we illustrate two versions of
the Bayes rule:  one (in red) based on the Kiefer-Wolfowitz estimate, the other
(in blue) based on the true mixing distributions.  The naive estimate appears as in
light grey line in this panel.  These Bayes rules are conditional
only on the observed $\bar y$ with the variance effect integrated out. Thus,
when we see a value of $\bar y$ near one of the mass points in $\{ -0.5, 1, 3\}$, 
the Bayes rule shrinks aggressively toward the corresponding $\mu$.
Between these values, the predicted $\mu$, being a conditional mean, takes 
intermediate values.  Extreme values of $\bar y$ in either tail again get
aggressively  shrunk toward the extreme points  of the estimated prior.

The situation we have just described is artificial in sense that we effectively
are assuming that we have observed $\bar y_i$ for each cross sectional unit,
but apparently have forgotten to compute the associated variance estimate.  If
we now rectify this oversight, we can consider a two dimensional Bayes rule
that maps $(\bar y_i , S_i)$ pairs into predictions of the $\mu_i$'s. Using
the same data and the estimates underlying Figure \ref{fig2.exv3} we illustrate
a contour plot of the this two dimensional Bayes rule in Figure \ref{fig2.exv3}.
We see that for central values of $\bar y_i$ the contours are essentially
vertical indicating the variance is uninformative about the mean, however
for outlying values of $\bar y_i$ the nonlinearity of the Bayes rule is apparent
with large observed variances making us more uncertain about the $\mu_i$'s.

When $\bar y$ is in the extremes,  the Bayes rule should shrink its estimate of $\mu$
to the extreme mass points at -0.5 and 3, but since the estimated prior has smaller
mass points nearby very extreme observations are attracted to these values.
In both tails one can see the effect of the variance estimate on this shrinkage effect;
when the estimated variance is small  then  there is more shrinkage to the nearest
mass point of the $\mu$ distribution.  When the observed variance is large, then
the posterior for $\mu$ is more evenly divided among several mass points and
consequently the posterior mean is more central.  For example,
when $\bar y = 1.5$ and the estimated variance is low, then
we can be quite confident that the observation comes from the $\mu = 1$ population.
Similarly when $\bar y = -1.5$  and the estimated variance is low, we can be
confident that this is a $\mu = -0.5$ observation.  However, in either of these
cases as the variance increases our confidence ebbs, and the Bayes rule
assigns more probability to the other nearby mass points.  For central values of
$\bar y$ the contours are nearly vertical indicating that the observed
variance is not informative in this region.  The observed pairs $(\bar y_i , S_i)$
are superimposed on the contours to give some sense of their dispersion.

\begin{figure}[ht]
  \begin{center}
          \resizebox{.8 \textwidth}{!}{{\includegraphics{Tweedie/fig2exv3.pdf}}}
  \end{center}
  \caption{{Bayes Rule for Gaussian location-scale mixture:  
  Based on the observations from Figure \ref{fig1.exv3} we
  illustrate the two dimensional Bayes rule for the mean parameter $\mu$ as
  a contour plot.}}
  \label{fig2.exv3}
\end{figure}

This form of the Bayes rule clearly illustrates that variances are informative
about the means in such circumstances, but it remains to be seen to what degree
they are quantitatively important in reducing the mean squared error of our
predictions of the mean.  To evaluate this we conducted a small simulation
experiment in which we generated 1000 replications from the model described
above and compared MSE's for 5 $\hat \mu_i$'s:  the naive $\hat \mu_i = \bar y_i$,
the Bayes Rule based only on the estimated distribution of means and variances
conditional on only the $\bar y_i$'s as in Figures \ref{fig1.exv3}, and conditional
on both $\bar y_i$ and $S_i$ as in Figure \ref{fig2.exv3}, and the two corresponding
Bayes rules based on the true mixing distributions.  The MSEs for these estimators
were 0.2800, 0.1575, 0.1595, 0.1570, and 0.1569 respectively.  Clearly, Bayes 
shrinkage is very successful in improving upon the naive predictions, and the
(Kiefer-Wolfowitz) empirical Bayes predictions are almost as good as
knowing the true mixing distribution, even for this modest, $n = 800, \; m = 11$ 
sample size, but the expected improvement from conditioning on the observed, 
individual-specific variance estimates is practically negligible.  Greater separation of the 
two variances might be expected to reveal more improvement.  But rather than 
pursuing further simulation exercises we now turn to 
a more realistic model including covariate effects.

\subsection{Covariate Effects}
Having seen how to estimate the Gaussian location-scale mixture model we will now
briefly describe how to introduce covariate effects into the model, which now takes the form,
\[
y_{it} = x_{it} \beta + \alpha_{i} + \sigma_{i} u_{it}.
\]
Given a $\beta$ it is easy to see that,
\[
\bar y_{i} | \alpha_{i}, \beta, \sigma_{i}^{2} \sim 
\mathcal{N}(\alpha_{i} + \bar x_{i} \beta, \sigma_{i}^{2})
\]
so the sufficient statistic for $\alpha_{i}$ is $\bar y_{i} - \bar x_{i} \beta$. 
Similarly, the sufficient statistic for $\sigma_{i}^{2}$ can be defined as,  
\[
S_{i} = \frac{1}{m_{i}-1} \sum_{t=1}^{m_{i}} ( y_{it} - x_{it} \beta - (\bar y_{i} - \bar x_{i} \beta))^{2}
\]
and $S_{i} | \beta, \sigma_{i}^{2} \sim \gamma (r_{i}, \sigma_{i}^{2}/r_{i})$, 
where as before, $r_{i} = (m_{i}-1)/2$. Apparently, using the familiar panel data terminology, 
the sufficient statistic for $\alpha_{i}$ contains the between information, while the 
within information, deviations from the individual means, is contained in the $S_{i}$. 
A note of caution should be added however since the orthogonality of the within and between 
information enjoyed by the classical Gaussian panel data model no longer holds in this 
general mixture setting. This can be seen more clearly by examining the likelihood function, 
\begin{align*}
L(\beta, h) &= \prod _{i=1}^{n} g ((\alpha, \beta, \sigma) | y_{i1}, \dots, y_{im_{i}}) \\
& = \prod_{i=1}^{n} \int \int \prod_{t=1}^{m_{i}} \sigma_i^{-1}
\phi((y_{it} - x_{it} \beta - \alpha_{i})/\sigma_{i})
h(\alpha_{i}, \sigma_{i}) d\alpha_{i} d\sigma_{i}\\
& = K \prod_{i=1}^{n}S_{i} ^{1-r_{i}} \int \int \sigma_i^{-1} 
\phi((\bar y_{i} - \bar x_{i} \beta - \alpha_{i})/\sigma_{i}) 
\frac{e^{-R_{i}}R_{i}^{r_{i}}}{S_{i} \Gamma(r_{i})} h(\alpha_{i}, \sigma_{i}) d\alpha_{i} d\sigma_{i} 
\end{align*} 
where $R_{i} = r_{i} S_{i} / \sigma_{i}^{2}$ and $K = \prod_{i=1}^{n}
\Big ( \frac{\Gamma(r_{i})}{r_{i}^{r_{i}}} (1/\sqrt{2\pi})^{m_{i}-1} ) \Big)$. 

Even with the independent prior assumption,  
$h(\alpha_{i}, \sigma_{i}) = h_\alpha(\alpha_{i}) h_{\sigma}(\sigma_{i})$, 
the likelihood does not factor because the Gaussian piece depends on both $\alpha_{i}$ and 
$\sigma_{i}$. However, the fact that $S_{i}$, hence the Gamma piece of the likelihood, 
does not depend on $\alpha_{i}$ provides a convenient estimation strategy by using the Gamma 
mixture to estimate $h_{\sigma}$, and a Studentized version of the Gaussian mixture, 
$(\bar y_{i} - \bar x_{i} \beta - \alpha_{i})/\sqrt{S_{i}} \sim t_{m_{i}-1}$, 
for estimating $h_{\alpha}$. Including covariates adapts this estimation strategy:
Given a $\beta$ we estimate the two mixing distributions and then evaluate the full profile
likelihood.  We will illustrate this approach in the next section.
This approach is related to  recent work by \citeasnoun{Bonhomme} on grouped patterns of
heterogeneity in panel data, although the estimation methods employed there based
on clustering algorithms are quite different.
 
\section{Empirical Bayesball}
In this section we describe our experience with methods described above to 
predict U.S. major league baseball batting averages.  

\subsection{The Data}
From \citeasnoun{espn} we have collected monthly data on the number of at bats
and hits for all U.S major league baseball players from 2002-2011, as well as
an indicator of whether the player is a pitcher.  We have aggregated this data
into half seasons to produce an unbalanced panel, observations
on players with more than 10 at bats in any half season,  and players with no less 
than 3 half-seasons,  leaving 1072 players and a total of 10,575 observations.
For the final subsection on age effects we have also collected the birth year of each player
from the ESPN website.
\subsection{The Model}
Following \citeasnoun{brown.08}, we consider transformed batting averages
\[
y_{it} = \arcsin \left ( \sqrt{\frac{H_{it} + 0.25}{N_{it} + 0.5}} \right )
\]
where $H_{it}$ denotes the number of ``hits'' of player $i$ in period $t$
and $N_{it}$ denotes his number of ``at bats'' in this period.  We will
assume, that the $y_{it}$'s are 
Gaussian with means, $\alpha_i = \arcsin(\sqrt{p_i})$, 
where $p_i$ is the individual specific batting success probability and variances,
$\sigma_i^2 \nu_{it}^2 = \sigma_i^2 / (4 N_{it})$.  Given an unbalanced panel 
with $t = 1, \cdots , m_i$ for $n$ players, we can define the sufficient 
statistics:
\[
\hat \alpha_i = (\sum_{t=1}^{m_i} \nu_{it}^{-2})^{-1}
\sum_{t=1}^{m_i} y_{it}/ \nu_{it}^2 
\]
and 
\[
S_i = \frac{1}{m_i -1} \sum_{t=1}^{m_i} (y_{it} - \hat \alpha_i)^2 /\nu_{it}^2.
\]
with Gaussian and Gamma distributions, respectively:
$\hat \alpha_i | \alpha_i , \sigma_i^2 \sim 
\NN ( \alpha_i , \sigma_i^2 v_i^2 )$ 
and $S_i | \sigma_i^2 \sim \gamma (r_i , \sigma_i^2/r_i)$,  where we set
$v_i^2 =  (4 \sum_{t=1}^{m_i} N_{it})^{-1}$ and $r_i = (m_i - 1)/2$.

For the sake of clarity, we will first discuss the case without covariates, 
introducing them in the final subsection.
We will assume that players draw a $\alpha_i$ and a $\sigma_i$ at random
from a distribution with density $h(\alpha_i , \sigma_i )$. 
Sufficiency then implies that we can factor the likelihood of the sample,
as a function of $(\alpha , \sigma ) \in \RR^{2n}$,
\begin{align*}
L(h)  &= \prod_{i=1}^n g ( (\alpha , \sigma )  \; | \; 
y_{i1} , \cdots , y_{im_i} )\\ 
&= \prod_{i=1}^n \int \int \prod_{t=1}^{m_i} \phi (( y_{it} - \alpha_i )/ 
\sigma_i \nu_{it} ) / (\sigma_i \nu_{it}) h(\alpha_i , \sigma_i ) 
d\alpha_i d\sigma_i \\
&= K \prod_{i=1}^n \int \int 
\phi (( \hat \alpha_{i} - \alpha_i )/ 
\sigma_i v_{i} ) / (\sigma_i v_{i}) 
\frac{e^{-R_i} R_i^{r_i}}{S_i \Gamma(r_i)}
h(\alpha_i , \sigma_i ) d\alpha_i d\sigma_i 
\end{align*}
where $R_i = r_i S_i / \sigma_i^2$.

As we have already noted, 
it is convenient at this point to make the further assumption that the
mixing density $h$ factors into $h = h_\alpha h_\sigma$.  In this case,
since the likelihood contribution of the $S_i$ is independent of the 
$\alpha_i$'s, we can solve the resulting (Kiefer-Wolfowitz) maximum
likelihood problem by first estimating the mixing density, 
$h_\sigma$ and  then estimating $h_\alpha$.  For the independent prior
case this is especially convenient since the likelihood factors into
independent gamma and Student $t$ components.  Estimation of the general
dependent prior case is beyond the scope of the present paper, but we
intend to explore this in future work. Estimation of the general dependent 
prior case raises no inherent technical difficulty except that it involves 
a two dimensional gridding.

\subsection{Estimation Results}

In Figure \ref{fig.fuv} we depict the estimated mixing densities for the
means and variances, $\mu$'s and $\theta$'s, for the model described in
the previous subsection.  In the transformed scale of the $\mu$'s we see
one large peak, and several smaller ones, with slight upper and lower 
``foothills.''  For the variance parameter $\theta$ we see only one very
pronounced peak slightly above one, and one smaller peak below
the large one.  Recall that the variance scale is relative to the binomial
model,  which implies that variance is completely determined by the longitudinal
sample size for each player.
Thus, a single peak at exactly one would imply exact adherence to the binomial
model.  Instead, we see a modest over-dispersion effect from the large peak,
and a much smaller peak corresponding to players exhibiting under-dispersion.
For the latter, more consistent, players prediction of future performance 
would presumably be easier.  The $C(\alpha)$ test of homogeneity of the variance
parameter, described in Section 3, described in detail in Appendix B, 
rejects with $p$-value 0.023, so we can be
reasonably confident that the smaller mass point is not spurious.


\begin{figure}[ht]
  \begin{center}
          \resizebox{ \textwidth}{!}{{\includegraphics{baseball/fuv.pdf}}}
  \end{center}
  \caption{{Estimated Mixing Distributions for $\mu$ and $\theta$ based on
  2002-11 longitudinal data.  Note that these distributions corresponds to
  the transformed, approximately Gaussian, observations.  See Figure \ref{fig.binp}
  for the untransformed estimated batting average distribution implied by this
  $f(\mu)$ distribution.  The $f(\theta)$ distribution depicted in the right panel can be
  interpreted as an estimate of a mixture of under and over dispersion components
  of the observed variances; the (larger) mass point with $\theta > 1$ shows that
  most players exhibit overdispersion relative to the binomial model, while the
  (smaller) mass point with $\theta < 1$ corresponds to a group of players that
  are less variable more consistent than predicted by the binomial model.}}
  \label{fig.fuv}
\end{figure}

Transforming the estimated mixing density, or prior, for the $\mu$'s back to
the natural scale of batting averages yields the distribution shown  in left panel of 
Figure \ref{fig.binp}.  Again, we see a similar configuration of peaks, but now 
located at more familiar places on the $[0,1]$ interval -- at least from baseball perspective.
In the right panel of Figure \ref{fig.binp} we illustrate the estimated mixing density
from the Kiefer-Wolfowitz MLE of a binomial model based on the aggregated baseball data
over the period 2002-11.  This estimate has the advantage of simplicity, but it neglects
the potential for over or under dispersion of the data.
We now turn to the question of how well we can predict future performance based
on these estimates.
\begin{figure}[ht]
  \begin{center}
	  \resizebox{ \textwidth}{!}{{\includegraphics{baseball/binp.pdf}}}
  \end{center}
  \caption{{Estimated Mixing Distributions for Batting Averages: In the left panel
  we present the transformed density based on the Gaussian model described above.
  And in the right panel we show the estimated mixing distribution for the MLE of
  the binomial model based on the aggregated performance of the players from
  2002-11.}}
  \label{fig.binp}
\end{figure}

\subsection{Prediction}

In this subsection we compare predictions based on the estimated models introduced above,
and compare their performance with alternative nonparametric Bayesian methods based on Dirichlet
priors.  Using the data on all 1072 players from the 2002-2011 seasons we estimate
our Gaussian model for $h_\alpha$ and $h_\theta$.  We drop pitchers from the prediction
exercise since they generally have fewer at bats.
Of the remaining 403 players with more than 40 at bats
in 2012, 336 had 3 or more half seasons prior to 2012 and were therefore qualified subjects
for our out-of-sample prediction exercise.  

In both panels of Figure \ref{fig.prediction}
we plot observed 2012 batting averages of these players against their prior (cumulated)
batting averages for the 2002-2011 period.  In the right panel of the figure we plot the
predicted batting averages based on the location-scale Gaussian model after retransformation
back to the probability scale.  The jagged nature of these predictions is due to the fluctuations
in the number of ``at bats'' in the  prior period as well as conditioning on prior variability,
$S_i$, which both influence the amount of shrinkage that is applied.  The grey line in both
plots represents the (45 degree) naive prediction, so we see that shrinkage is greatest in the
lower tail while the estimated mass points of the prior distribution in the upper tail tend to
reduce the amount of shrinkage there.  In the left panel we have plotted the empirical Bayes
predictions from the somewhat simpler binomial version of the model.  The strong resemblance
of the two sets of predictions is indicative of the fact that the Gaussian transformation model
and its attempt to account for over and under dispersion of the binomial model did not alter
the fundamental nature of the binomial model results.  Clearly, the dispersion of the plotted
points in both panels starkly reflects the difficulty of the prediction exercise.  

\begin{figure}[ht]
  \begin{center}
          \resizebox{ \textwidth}{!}{{\includegraphics{profile/bayesrule_f8.pdf}}}
  \end{center}
  \caption{Prediction of 2012 Batting Averages:  In the left panel we plot the 
  Bayes rule prediction based on the Binomial model in contrast to the naive 
  prediction (grey line). In the right panel, the solid blue line is the empirical 
  Bayes rule prediction based on the Gaussian model without any covariates and 
  the dotted red line is the prediction taking into account of the quadratic 
  effects of age as discussed in the last subsection. The actual outcome of 
  2012 batting average is superimposed in both figures.}
  \label{fig.prediction}
\end{figure}

To compare the performance of the foregoing predictions with more conventional Bayesian methods
we have also considered an alternative nonparametric formulation of the binomial model employing
Dirichlet process priors as introduced by \citeasnoun{Ferguson73}, \citeasnoun{Antoniak74}
and \citeasnoun{Ferguson83}.  Deferring a discussion of the details to Appendix B, we consider
two versions of the Dirichlet model both with Dirichlet prior $\DD (\alpha G_0 )$ and
base distribution $G_0$ taken as $\mbox{Beta}(1,1)$, i.e. uniform on $[0,1]$.  
The scaling parameter, $\alpha$ is either 0.01 or 1, both
reflecting relative ignorance about the prior.  Markov chain Monte-Carlo methods produce a posterior
of the predictive density for the mixing distribution corresponding to the point estimate 
produced by the Kiefer-Wolfowitz estimator.  
The mean of this posterior serves as a predictive density for the mixing 
distribution of the binomial parameter can thus be used to produce predictions for our 
out-of-sample 2012 players as for the Binomial model described above.  

\begin{figure}[ht]
  \begin{center}
          \resizebox{ \textwidth}{!}{{\includegraphics{baseball/figDP.pdf}}}
  \end{center}
  \caption{{Dirichlet Estimates of the Mixing Distribution:  In the left panel we plot
  the estimated mixing distribution based on the Dirichlet prior with precision parameter
  $\alpha= 1$.  In the right panel prior precision has been reduced to
  $\alpha= 0.01$. These estimates may be interpreted as regularized versions of those
  produced by the Kiefer-Wolfowitz MLE appearing in Figure \ref{fig.binp}.}}
  \label{fig.DPmix}
\end{figure}

It will come as no surprise, given the appearance of Figure \ref{fig.prediction}, that
none of the methods perform very well, and given the similarity of the estimated mixing
distribution, performance of their predictions are also quite similar.
This is indeed the case: the root mean squared error for the Gaussian and Binomial models
are respectively, 0.0393 and 0.0395.  The corresponding root mean squared errors for
the two Dirichlet models are both 0.0395, and this is no better than the naive prediction
based on prior accumulated batting averages.  The explanation for this rather dismal 
performance is partly that the mixture models estimated multimodality doesn't
generate much shrinkage, and ultimately even prior performance  isn't a very good
gauge of future batting success.

\subsection{Age and the Skill of Batting}

\begin{figure}[ht]
  \begin{center}
	  \resizebox{ .8\textwidth}{!}{{\includegraphics{profile/profage.pdf}}}
  \end{center}
  \caption{Profile Likelihood for a Linear Age Effect on Batting Average:
  The Wilks confidence interval for the age effect is represented by the shaded region.}
  \label{fig.profage}
\end{figure}

\begin{figure}[ht]
  \begin{center}
	  \resizebox{ .8\textwidth}{!}{{\includegraphics{profile/profage2.pdf}}}
  \end{center}
  \caption{Contours of  Profile Likelihood for a Quadratic Age Effect on Batting Average:
  The Wilks confidence region for the age effect parameters is represented by the shaded region.}
  \label{fig.profage2}
\end{figure}

\begin{figure}[ht]
  \begin{center}
	  \resizebox{ .8\textwidth}{!}{{\includegraphics{profile/ageeff.pdf}}}
  \end{center}
  \caption{The Estimated Quadratic Age Effect:  The vertical axis is in units of the transformed
  batting average.}
  \label{fig.ageeff}
\end{figure}

Finally, to illustrate the role that covariates can play in the empirical Bayes approach
we have developed, we consider a model in which player's age influences his batting average.
Following the approach outlined in the previous section we compute the profile likelihood
on a grid of $\beta$ parameters.  The only required modification is the reweighting by the
$v_{it}$'s.  When we include only a linear age effect, we obtain the profile likelihood
shown in Figure \ref{fig.profage} that yields a point estimate of about $\beta = -0.12$.
The usual Wilks confidence interval also illustrated in the figure is $[-0.151, -0.092]$.  This relentless
deterioration in batting skill seems a bit implausible, and depressing, so we have also
tried a simple quadratic alternative.  Now a contour plot of the profile likelihood is
shown as Figure \ref{fig.profage2} and we find, after unrescaling,  that the age effect appears
as shown in Figure \ref{fig.ageeff}.  The peak occurs at 26.5 years, by age 37 performance
has declined by about 0.02 units in the transformed batting average, which corresponds to a
decline of a typical batting average of 0.320 at the peak to 0.302.  The well-known baseball 
guru Bill James is on record as asserting that batting ability peaks at age 27, so our estimates
are at least consistent with his casual empiricism.  Our estimated age profile also bears an
uncanny resemblance to the age profile of Ty Cobb (1905-1928) analyzed in \citeasnoun{morris}. Incorporating the quadratic age effects improves the prediction performance 
slightly, so that the root mean squared error reduces to 0.0378. The Bayes rule 
prediction is plotted as the dotted line in Figure \ref{fig.prediction}.

\section{Conclusion}

Models of unobserved heterogeneity for longitudinal data are common in applied econometrics.
We have argued that empirical Bayes methods based on nonparametric maximum likelihood estimation 
of mixture models offer a natural formulation of these models.  Recent developments in convex
optimization greatly facilitate estimation of such models.  Semiparametric versions of these
models including covariate effects are shown to be effectively analyzed with profile likelihood.  
A potential criticism of the foregoing approach is that it requires us to assume a parametric
form for the base distribution, in our setting the Gaussian.  Of course, location-scale mixtures
of Gaussians is quite a general class, so from a prediction perspective the normality assumption
is not especially onerous.  Moreover, in our baseball application normality is easily justified
since the transformed batting averages have a strong claim to approximate normality.  In some
special cases it may be possible to {\it estimate} the base distribution nonparametrically from
extraneous sample information.  Recent work of \citeasnoun{BS} illustrates an educational treatment effect
application of this type for which a deconvolution strategy based on empirical characteristic functions
is employed.   We hope to explore the extension of Kiefer-Wolfowitz mixture methods to such settings
in future work.

\bibliography{brown,MCMC}
\appendix
{\footnotesize
\section{Tweedie's Formulae (Proof of Propositions 1 and 2)} 
The simplest way to derive Tweedie's Formula for the Gaussian case seems to 
be to consider the more general exponential family compound decision problem 
in which
\[
g (y) = \int \varphi(y , \theta) dF(\theta),
\]
where $\varphi$ is a known exponential family density with natural parameter
$\theta$, so we may write,
\[
\varphi (y , \theta ) = m(y) e^{y \theta} h(\theta),
\]
and $F$ is again a mixing distribution over the parameter $\theta$.
Quadratic loss implies that the Bayes rule is the conditional mean:
\begin{eqnarray*}
\delta_G (y) & = & \EB [ \Theta | Y = y]\\
&=& \int \theta \varphi(y,\theta) dF / \int  \varphi(y,\theta) dF\\
&=& \int \theta e^{y \theta} h(\theta) dF / \int e^{y \theta} h(\theta) dF \\
&=& \frac{d}{dy} \log ( \int  e^{y \theta} h(\theta) dF \\
&=& \frac{d}{dy} \log ( g(y) /m(y)) 
\end{eqnarray*}
Differentiating again,
\begin{align*}
\delta^\prime (y) = \frac{d}{dy} 
\left [ \frac { \int \theta \varphi dF } { \int \varphi dF } \right ]
& =  \frac { \int \theta^2 \varphi dF } { \int \varphi dF } 
 - \left ( \frac { \int \theta \varphi dF } { \int \theta \varphi dF } \right )^2 \\
& =  \EB [ \Theta^2 | Y = y] - (\EB [ \Theta | Y = y])^2\\
& =  \VB [ \Theta | Y = y] \geq 0,
\end{align*}
implying that $\delta$ must be monotone.  
When $\varphi$ is Gaussian we have,
\[
\varphi (y, \theta) = \phi(y - \theta) = K \exp \{ - (y - \theta)^2 / 2 \}
= K  e^{-y^2/2} \cdot e^{y \theta} \cdot e^{- \theta^2 /2}, 
\]
so $m(y) = e^{-y^2/2} $ and the logarithmic derivative yields our Bayes rule.

For the gamma mixture case there is an analogous formula, for the 
natural parameter $-r/\theta$, proceeding as before, 
\[
\tilde{\delta}(x)=E[-r/\theta | X=x]
= \frac{\int -\frac{r}{\theta}f(x|\theta)dG(\theta)} {\int f(x|\theta)dG(\theta)}
=\frac{d}{dx}log(\frac{g(x)}{x^{r-1}})
=\frac{g'(x)}{g(x)}-\frac{r-1}{x}.
\]
If, on the other hand, we would, quite naturally, like to compute the
expectation of the {\it unnatural} parameter $\theta$, then we obtain instead,
\begin{align*}
\delta (x) & =  \EB [ \Theta | X = x]\\
&= \int \theta \gamma(x,\theta) dF (\theta) / \int  \gamma(x,\theta) dF (\theta)\\
&= \int \frac{\theta}{\Gamma(r) (\theta /r)^r} x^{r-1} 
\exp( - x r / \theta ) dF (\theta) / g(x) \\
&= r x^{r-1} \int \frac{(r/\theta)^r}{\Gamma(r)} \frac{\theta}{r}) 
\exp( - x r / \theta ) dF (\theta) / g(x) \\
&= r x^{r-1} \int \frac{(r/\theta)^r}{\Gamma(r)} \int_x^\infty 
\exp( - x r / \theta ) dF (\theta) / g(x) \\
&= r x^{r-1} \int_x^\infty y^{r-1} \int \gamma(y|\theta) dF (\theta) dy /g(x) \\
&= r x^{r-1} \int_{x}^{\infty}y^{1-r}g(y)dy /g(x).
\end{align*}
It remains to show that this formulation of the Bayes rule is monotone:
\begin{align*}
\delta^\prime (x)&=\frac{r(r-1)x^{r-2}\int_{x}^{\infty}y^{1-r}g(y)dy+rx^{r-1}(-x^{1-r})g(x)}{g(x)}\\
& \qquad \qquad -\frac{rx^{r-1}\int_{x}^{\infty}y^{1-r}g(y)dy g'(x)}{g^{2}(x)}\\
&=\frac{r x^{r-1}\int_{x}^{\infty}y^{1-r}g(y)dy}{g(x)}\left [\frac{r-1}{x}-\frac{g'(x)}{g(x)}\right ]-r\\
&=- \delta(x)\tilde{\delta}(x)-r\\
&=\delta(x)\left [E[\frac{r}{\theta}|X=x]-\frac{r}{E[\theta|X=x]}\right ]\\
&=\delta(x) r\left [E[\frac{1}{\theta}|X=x]-\frac{1}{E[\theta|X=x]}\right ],\\
\end{align*}
which is positive by Jensen's inequality.


\section{ $C(\alpha)$ test for Gamma mixture}
We provide details here for the $C(\alpha)$ strategy for testing homogeneity of the scale 
parameter in the Gaussian location-scale mixture model considered in Section 2.2. 
Given the sufficient statistics $S_{i}$, we would like to test whether they arose from a
Gamma distribution with $\sigma_{i}^{2} = \sigma_{0}^{2}$ for all $i$, or equivalently 
that $F_{\sigma^{2}}$ is a degenerate distribution with mass one at $\sigma_{0}^{2}$. 
We will reparameterize by writing $\sigma_{i}^{2}= \sigma_{0}^{2} \exp (\xi U_{i})$ where 
$U_{i} \sim H$ are iid with $\mathbb{E}(U)=0$ and $V(U)=1$. The homogeneity test is 
therefore: $H_{0}: \xi =0$ vs. $H_{1}: \xi  \neq 0$. \\

Following the $C(\alpha)$ procedure for mixture models as in \citeasnoun{Gu13}, 
the score for $\xi$ and the nuisance parameter $\sigma_{0}^{2}$ are respectively:
\[
\begin{array}{ll}
v_{1i} & := \nabla_{\xi}^{2} \log f_{i} |_{\xi=0} = (s_{i} r_{i}/\sigma_{0}^{2} - r_{i})^{2} - s_{i} r_{i} / \sigma_{0}^{2}\\
v_{2i} & := \nabla_{\sigma_{0}^{2}} \log f_{i} | _{\xi =0} = s_{i}r_{i}/\sigma_{0}^{4} - r_{i}/\sigma_{0}^{2}
\end{array}
\]
where $f_{i}$ refers to the Gamma mixture density for the $S_{i}$, and $r_i = (m_i - 1)/2 $.
The efficient score for $\xi$ is therefore, 
\[
g_{i}(\sigma_{0}^{2})  = v_{1i}- \frac{cov(v_{1}, v_{2})}{V(v_{2})} v_{2i}  
= (Y_{i} - r_{i})^{2} - Y_{i} - (Y_{i}-r_{i})
\]
where $Y_{i} = S_{i}r_{i}/\sigma_{0}^{2}$ with moments 
$\mathbb{E}(Y_{i}^{k})= \prod _{j=0}^{k-1} (r_{i} +j)$. 
Replacing the nuisance parameter by its MLE 
$\hat \sigma_{0}^{2} = \sum_{i} s_{i} r_{i} / \sum_{i} r_{i}$, the $C(\alpha)$ test statistic is,
\[
T_{n}  =  \frac{\sum_{i} g_{i}(\hat \sigma_{0}^{2})}{\sqrt{\sum_{i} V(g_{i}(\hat \sigma_{0}^{2}))}} = \frac{\sum_{i} \Big [ (s_{i}r_{i}/\hat \sigma_{0}^{2} - r_{i})^{2} - s_{i} r_{i} /\hat \sigma_{0}^{2} \Big ] }{\sqrt{ \sum_{i} (2 r_{i}^{2} + 2 r_{i})}}
\]
Under $H_{0}$, $T_{n} \sim \mathcal{N}(0,1)$, 
and we reject the null at significance level $\alpha$ if $T_{n} >  \Phi^{-1}(1-\alpha)$. 
\section{Dirichlet Process Methods for Compound Decisions}


\citeasnoun{Ferguson73} first proposed using the Dirichlet process as a prior distribution 
for estimating an unknown probability measure $P$. 
\citeasnoun{Antoniak74} and \citeasnoun{Ferguson83} extend the Dirichlet process to a 
mixture of Dirichlet processes, suited to compound decision problems like the binomial mixture 
model we have considered above. We do not observe samples on $\mu_{1}, \dots, \mu_{n}$ directly, 
instead we would like to find the posterior distribution of $\mu$ conditional on observing 
$y_{1}, \dots, y_{n}$. \citeasnoun{Antoniak74} shows that if $\mu_{1}, \dots, \mu_{n}$ are 
drawn from a distribution $F$, with Dirichlet process prior $\mathcal{D}(\alpha G_{0})$, 
and if $Y_{1}, \dots, Y_{n}$ are i.i.d. random variable with density 
$g(y)=\int f(y | \mu) dF(\mu)$, then the posterior distribution of $F$ given 
$y_{1}, \dots, y_{n}$ is a mixture of Dirichlet processes:
\[
F \mid y_{1}, \dots, y_{n} \sim \int \dots \int \mathcal{D}(\alpha G_{0} + \sum_{i=1}^{n} \delta_{\mu_{i}}) dH(\mu_{1}, \dots, \mu_{n} | y_{1}, \dots, y_{n})
\]

The above posterior seems rather intractible, however with the help of the P\'olya urn 
representation of the Dirichlet process by \citeasnoun{Blackwell73}, one obtains a simplified 
form of the posterior distribution, (see, for example, \citeasnoun{Lo84}). We will see that 
this representation opens the way for MCMC sampling from the posterior distribution.

The P\'olya urn representation provides the following conditional distribution, 
\begin{equation} \label{eq1}
\mu_{i} | \mu_{1}, \dots, \mu_{i-1} \sim \frac{\alpha}{\alpha +i-1}G_{0}(\mu_{i}) + \frac{1}{\alpha+i-1} \sum_{j=1}^{i-1} \delta_{\mu_{j}}(\mu_{i})
\end{equation}
where $\delta_{\mu_{j}}$ denotes the (Dirac) distribution with mass one at point $\mu_{j}$. 
In the P\'olya urn scheme each new draw $\mu_{i}$ is a random draw from the 
base distribution $G_{0}$ with probability $\frac{\alpha}{\alpha+i-1}$ and 
otherwise take the value of the previous $i-1$ $\mu$'s, each with probability 
$1/(\alpha+i-1)$. The joint distribution is therefore,  
\[
p(\mu_{1}, \dots, \mu_{n})= \prod_{i=1}^{n} 
\frac{\alpha G_{0}(\mu_{i})+\sum_{j=1}^{i-1}\delta_{\mu_{j}}(\mu_{i})}{\alpha+i-1},
\]
and the posterior distribution of $(\mu_{1}, \dots, \mu_{n} | y)$ can thus be written as,
\begin{equation} \label{post}
p(\mu_{1}, \dots, \mu_{n} |y) \propto \prod_{i=1}^{n} f(y_{i} | \mu_{i})  \frac{\alpha G_{0}(\mu_{i}) + \sum_{j=1}^{i-1}\delta_{\mu_{j}}(\mu_{i})}{\alpha + i-1}. 
\end{equation}
Note that we have integrated out the infinite dimensional object $F$, the distribution 
for  $\mu$,  

\subsection{Gibbs Sampling I}
Our aim is to sample from the above posterior distribution. This can be accomplished by 
simulating a Markov chain that has the posterior as its stationary equilibrium distribution. 
The Gibbs sampling algorithm in \citeasnoun{EscobarWest94} provides a simple way to do this. 
It relies on the simple form of the posterior distribution of $\mu_{i}$ conditional on the 
remaining parameters $\mu_{(-i)} = (\mu_{1}, \dots, \mu_{i-1}, \mu_{i+1}, \dots, \mu_{n})$ 
and data $\bf{y}$ such that we only need to track one variable at each step instead of all $n$ of them.
The easiest implementation chooses $G_{0}$ as the conjugate prior. 
For the binomial setting, this leads us to choose $G_{0}$ as a beta distribution with 
density $\mathcal{B}(a,b)$. 
%A normal mixture example is provided in \citeasnoun{Escobar94} and \citeasnoun{EscobarWest94}. \\

We can write the posterior distribution as, 
\begin{align} \label{eq2}
p(\mu_{i} | \mu_{(-i)}, \bf{y}) & = \frac{ f(y_{i} | \mu_{1}, \dots, \mu_{n}) \cdot p( \mu_{i} | \mu_{(-i)}) }{\int f(y_{i} | \mu_{1}, \dots, \mu_{n}) \cdot p( \mu_{i} | \mu_{(-i)})  d\mu_{i}} \nonumber \\ 
& \propto  \frac{\alpha}{\alpha+n-1} \mathcal{B}(a,b) f(y_{i} | \mu_{i}) + \frac{1}{\alpha +n-1} \underset{j \neq i}{\sum} f(y_{i} | \mu_{j} ) \delta_{\mu_{j}} \nonumber \\
& \propto  \frac{\alpha}{\alpha + n-1} \frac{B(y_{i}+a, l_{i}-y_{i}+b)}{B(a, b)} \mathcal{B}(y_{i}+a, l_{i}-y_{i}+b)  \nonumber \\
& + \frac{1}{\alpha + n -1} \underset{j \neq i}{\sum} \mu_{j} ^{y_{i}}(1-\mu_{j})^{l_{i}-y_{i}} \delta_{\mu_{j}}
\end{align}
%with the normalizing constant as $(\frac{\alpha}{\alpha+n-1} 
%\frac{B(y_{i} +a , l_{i}-y_{i}+b)}{B(a,b)} + 
%\frac{1}{\alpha +n-1} \underset{j \neq i}
%{\sum} \mu_{j}^{y_{i}}(1-\mu_{j})^{l_{i}-y_{i}})^{-1}$ 
where $B(\cdot, \cdot)$ denotes the Beta function, $\mathcal{B}(\cdot, \cdot)$ denotes the 
beta density and $f(y_{i}|\mu_{i})$ is the binomial density. Note that the posterior 
distribution only depends on $y_{i}$ because for $j \neq i$, $y_{j}$ is conditionally 
independent of $\mu_{i}$ given $\mu_{j}$'s. 

Gibbs sampling initiates the Markov chain by sampling $(\mu_{1}^{(0)}, \dots, \mu_{n}^{(0)})$ 
from $\mathcal{B}(y_{i}+a, l_{i}-y_{i}+b)$ (the posterior obtained by updating the prior 
$G_{0}$ via the likelihood $f(y_{i} | \mu_{i})$ up to a normalization constant). 
The first step of the Markov Chain is defined as: 

Sample $\mu_{1}^{(1)}$ from $p(\mu_{1}| \mu_{2}^{(0)}, \dots, \mu_{n}^{(0)}, \bf{y})$. 

Sample $\mu_{2}^{(1)}$ from $p(\mu_{2} | \mu_{1}^{(1)}, \mu_{3}^{(0)}, \dots, \mu_{n}^{(0)}, 
\bf{y})$.\\ $\dots$

Sample $\mu_{n}^{(1)}$ from $p(\mu_{n} | \mu_{1}^{(1)}, \dots, \mu_{n-1}^{(1)}, \bf{y})$. 

Continuing the iteration, the chain stabilizes at its equilibrium distribution and the 
resulting sample constitute draws from the distribution with density function 
$p(\mu_{1}, \dots, \mu_{n} | \bf{y})$ as in (\ref{post}). 

\subsection{Gibbs Sampling II}
Gibbs sampling as above may have rather slow convergence because of the clustering behavior of 
the $\mu_{i}$'s (as already noted by \citeasnoun{Antoniak74}). As discussed in \citeasnoun{Neal}, 
since the Gibbs sampling algorithm implemented above can not change $\mu$ for more than one 
observation at a time, changes of $\mu$ values for observations in the same cluster occur 
rather rarely, leading to rigidity in the chain and hence slow convergence. 
This can be avoided by introducing a latent class variable as implemented in 
\citeasnoun{Bush} and \citeasnoun{West94}. It is also the approach adopted by the 
DPbetabinom function in the  R package DPpackage that we employ.

The latent class model is equivalent to the model introduced above, except that we 
specify a configuration variable $\mathcal{S}=(\mathcal{S}_{1}, \dots, \mathcal{S}_{n})$ 
that classifies the $\mu_{i}$'s into $k$ distinct clusters. 
Denoting $n_{j}= \# \{ \mathcal{S}_{i} =j \}$ for $j = 1, \dots, k$, the distinct values of 
$\mu$ form the set $\{ \theta_{1}, \dots, \theta_{k} \}$ and we set
$I_{j} = \{ i: \mathcal{S}_{i}=j, i=1, \dots, n\}$. 
As discussed in \citeasnoun{Antoniak74} and \citeasnoun{West90}, the $\theta_{i}$ are 
random samples from $G_{0}$ and $k$ is related to the sample size $n$ and the precision 
parameter $\alpha$. Given $k$, the $\mu_{i}$'s are selected from the set of $\theta$ 
according to a multinomial distribution. With the above described structure, we have 
the distribution for $\mu_{i}$ conditional on $\mu_{(-i)}$ as, 
\[
p(\mu_{i} \mid \mu_{(-i)}, \mathcal{S}_{(-i)},  k_{(-i)}) \sim \frac{\alpha}{\alpha+n-1} G_{0} + \frac{1}{\alpha+n-1} \sum_{j=1}^{k_{(-i)}}  n_{j}^{(-i)} \delta_{\theta_{j}^{(-i)}}
\]
where $\mathcal{S}_{(-i)}$ is the configuration corresponding to all entries in 
$\mu_{(-i)}$ and $k_{(-i)}$ is the number of distinctive values of $\mu_{(-i)}$, 
contained in the set $\theta^{(-i)}$ and $n_{j}^{(-i)} = \# \{ \mathcal{S}_{(-i)} = j \}$ 
for $j = 1, \dots, k_{(-i)}$. We still have a P\'olya urn scheme:  with probability 
$\alpha/(\alpha+n-1)$, $\mu_{i}$ is a random draw from $G_{0}$, and otherwise takes a value 
from the set $\theta^{(-i)}$ with probability proportional to the multinomial counts 
$n_{j}^{(-i)}$. 

Choosing $G_{0}$ to be the conjugate prior $\mathcal{B}(a,b)$, the posterior 
distribution of $\mu_{i}$ conditional on data $\bf{y}$ and $\mu_{(-i)}$ is thus, 
\[
\begin{array}{ll}
p(\mu_{i} \mid \mu_{(-i)}, \mathcal{S}_{(-i)},  k_{(-i)}, y)  & \propto \frac{\alpha}{\alpha+n-1} f(y_{i} | \mu_{i}) \mathcal{B}(a,b) + \frac{1}{\alpha +n-1} \sum_{j=1}^{ k_{(-i)}}  n_{j}^{(-i)} f(y_{i} | \theta_{j}^{(-i)}) \delta_{\theta_{j}^{(-i)}}\\
& = q_{i,0} \mathcal{B}(y_{i}+a, l_{i}-y_{i}+b) + \sum_{j=1}^{k_{(-i)}} q_{i,j} \delta_{\theta_{j}^{(-i)}}
\end{array}
\]
with 
\[
q_{i,j}=
\begin{cases}
c \frac{\alpha}{\alpha+n-1} \frac{B(y_{i}+a, l_{i}-y_{i}+b)}{B(a,b)} & \mbox{  if   }   j=0\\
c \frac{n_{j}^{(-i)}}{\alpha+n-1} (\theta_{j}^{(-i)})^{y_{i}}(1-\theta_{j}^{(-i)})^{l_{i}-y_{i}} &\mbox{  if   } j >0
\end{cases}
\]
and normalizating constant,
\[
 c^{-1} = \frac{\alpha}{\alpha+n-1} \frac{B(y_{i}+a, l_{i}-y_{i}+b)}{B(a,b)}+ \frac{1}{\alpha+n-1} \sum_{j=1}^{k_{(-i)}} n_{i}^{(-i)} (\theta_{j}^{(-i)})^{y_{i}}(1-\theta_{j}^{(-i)})^{l_{i}-y_{i}}.
 \]
 
This second Gibbs sampling algorithm differs from the previous one in that instead of 
iteratively updating $\mu_{i}$'s, we update the configuration variable $\mathcal{S}$ 
according to its posterior distribution, that is, 
\[
P(\mathcal{S}_{i} =j \mid y, \mu_{(-i)}, \mathcal{S}_{(-i)}, k_{(-i)}) = q_{i,j}.
\]
We initiate the chain by choosing $\mathcal{S}_{i}^{(0)}=i$ (i.e., each observation forms 
its own cluster) and $\theta^{(0)}$'s can be drawn from $\mathcal{B}(y_{i}+a, l_{i}-y_{i}+b)$. 
The Markov chain is simulated as,
\begin{enumerate}
\item Given values for $\theta^{(0)}$ and $\mathcal{S}^{(0)}$, generate a new 
	$\mathcal{S}^{(1)}$ according to posterior distribution specified above successively. 
	For any index $i$ that has $\mathcal{S}_{i}^{(1)}=0$, draw a new $\mu_{i}$ 
	from $\mathcal{B}(y_{i}+a, l_{i}-y_{i}+b)$. Count the number of clusters $k^{(1)}$.

\item Given $k^{(1)}$ and $\mathcal{S}^{(1)}$, generate a new set of $\theta^{(1)}$ 
	by sampling from 
\[
p(\theta_{j} | y, \mathcal{S}^{(1)}, k^{(1)}) \propto \prod_{r \in I_{j}^{(1)}} f(y_{r}|\theta_{j})dG_{0}(\theta_{j})
\]

\item  Continue iterating \dots
\end{enumerate}

After discarding the first few steps of the Markov chain (burn-in), we can use the 
simulated sample for the following posterior analysis. 





\subsection{Posterior Analysis}
Given the posterior distribution (\ref{eq2}), the posterior mean for $\mu_{i}$ 
from the $m^{th}$ step of the MCMC scan is 
\begin{align} \label{pmean}
\mathcal{E}_{i}^{(m)}=\mathbb{E}(\mu_{i}^{(m)} | y, \mu_{(-i)}^{(m)}) & = \int \mu_{i} p(\mu_{i} | \mu_{(-i)}, y) d\mu_{i} \nonumber \\
&= \frac{\alpha \frac{B(y_{i} +a +1, l_{i}-y_{i} +b)}{B(a,b)} + \sum_{j \neq i} (\mu_{j}^{(m)}) ^{y_{i}+1} ( 1- \mu_{j}^{(m)})^{l_{i}-y_{i}}}{\alpha \frac{B(y_{i}+a, l_{i}-y_{i}+b)}{B(a,b)} + \sum_{j \neq i} (\mu_{j}^{(m)})^{y_{i}}(1-\mu_{j}^{(m)})^{l_{i}-y_{i}}} 
\end{align} 

The posterior mean of $\mu_{i}$ given data using the entire chain is thus estimated by, 
\[
\frac{1}{M}\sum_{m=1}^{M} \mathcal{E}_{i}^{(m)},
\]
where $M$ is the total number of MCMC scans after initial burn-in. 
As noted by \cite{GelfandSmith90} this is essentially Rao-Blackwellization. 
The posterior variance can be found accordingly as,  
\[
\mathbb{V}(\mu_{i} | y) = \mathbb{E}( \mathbb{V}(\mu_{i} | \mu^{(-i)}, y)) + \mathbb{V} ( \mathbb{E}(\mu_{i} | \mu^{(-i)}, y)). 
\]
\\
\subsection{Predictive distribution}
Nonparametric Bayesian analysis often focuses on the predictive distribution of $\mu_{n+1}$ or 
$y_{n+1}$ for a future experiment.  It is also something routinely reported by software packages.
%\footnote{For example, 
%``densm.p'' in DPbetabinom function of the DPpackage is the predictive density for $\mu_{i}$ 
%evaluated at grid points on the $[0,1]$ domain.}   
As noted by \citeasnoun{Liu96}, since $F$ is an infinite-dimensional object, there is no 
easy way of exploring its full posterior distribution. However, we can look at its 
posterior mean, $\mathbb{E}(F | y )$, which also turns out to be the predictive distribution 
of a future $\mu_{n+1}$. In particular, given the latent class model , the predictive density 
for a future $\mu$ in the $m^{th}$ round of Markov Chain simulation, evaluated at grid points 
on the domain of $(0,1)$, is,
\[
\begin{array}{ll}
p(\mu_{n+1}^{(m)} \mid y, \mathcal{S}^{(m)}, k^{(m)}) &= \mathbb{E}(F \mid y_{1}, \dots, y_{n}) \\
&= \frac{\alpha}{\alpha+n} G_{0}(\cdot)+ \frac{1}{\alpha +n} \sum_{j=1}^{k^{(m)}} n_{j}^{(m)} \mathcal{B}(\cdot; Z_{j}^{(m)} +a, L_{j}^{(m)}-Z_{j}^{(m)}+b)
\end{array}
\]
with $Z_{j}^{(m)}=\sum_{r \in I_{j}^{(m)}}y_{r}$ and 
$L_{j}^{(m)}=\sum_{r \in I_{j}^{(m)}}l_{r}$. 
The predictive distribution is then obtained  as,
\[
\frac{1}{M}\sum_{m=1}^{M} p(\mu_{n+1}^{(m)} \mid y, \mathcal{S}^{(m)}, k^{(m)}).
\]



Given this predictive density we can compute the Bayes rule for predicting $\mu$ given an
observed $y$ as in the binomial Kiefer Wolfowitz setting,
\[
\mathbb{E}(\mu \mid y) =\frac{\int \mu f(y \mid \mu) \hat f(\mu) d\mu}{\int f(y \mid \mu) \hat f(\mu) d\mu}
\]
The predictive density depends on the Dirichlet prior $\mathcal{D}(\alpha G_{0})$. 
In the application, we take the distribution $G_{0}$ to be $\mathcal{B}(1,1)$, 
i.e., the uniform distribution on the interval $[0,1]$. The precision parameter 
$\alpha$ takes values $1$ and $0.01$. 
The bigger $\alpha$ is, the more confidence we have in $G_{0}$. 
}
\end{document}


